%=============================================================================
% Introduction and Notation
%=============================================================================
\chapter*{Introduction}
\addcontentsline{toc}{chapter}{Introduction}

Commutative algebra is essentially the study of commutative rings. Roughly speaking, it has developed from two sources: (1) algebraic geometry and (2) algebraic number theory. In (1) the prototype of the rings studied is the ring $k[x_1, \ldots, x_n]$ of polynomials in several variables over a field $k$; in (2) it is the ring $\Z$ of rational integers. Of these two the algebro-geometric case is the more far-reaching and, in its modern development by Grothendieck, it embraces much of algebraic number theory. Commutative algebra is now one of the foundation stones of this new algebraic geometry. It provides the complete local tools for the subject in much the same way as differential analysis provides the tools for differential geometry.

This book grew out of a course of lectures given to third year undergraduates at Oxford University and it has the modest aim of providing a rapid introduction to the subject. It is designed to be read by students who have had a first elementary course in general algebra. On the other hand, it is not intended as a substitute for the more voluminous tracts on commutative algebra such as Zariski-Samuel \cite{ZariskiSamuel} or Bourbaki \cite{Bourbaki}. We have concentrated on certain central topics, and large areas, such as field theory, are not touched. In content we cover rather more ground than Northcott \cite{Northcott} and our treatment is substantially different in that, following the modern trend, we put more emphasis on modules and localization.

The central notion in commutative algebra is that of a prime ideal. This provides a common generalization of the primes of arithmetic and the points of geometry. The geometric notion of concentrating attention ``near a point'' has as its algebraic analogue the important process of localizing a ring at a prime ideal. It is not surprising, therefore, that results about localization can usefully be thought of in geometric terms. This is done methodically in Grothendieck's theory of schemes and, partly as an introduction to Grothendieck's work \cite{Grothendieck}, and partly because of the geometric insight it provides, we have added schematic versions of many results in the form of exercises and remarks.

The lecture-note origin of this book accounts for the rather terse style, with little general padding, and for the condensed account of many proofs. We have resisted the temptation to expand it in the hope that the brevity of our presentation will make clearer the mathematical structure of what is by now an elegant and attractive theory. Our philosophy has been to build up to the main theorems in a succession of simple steps and to omit routine verifications.

Anyone writing now on commutative algebra faces a dilemma in connection with homological algebra, which plays such an important part in modern developments. A proper treatment of homological algebra is impossible within the confines of a small book: on the other hand, it is hardly sensible to ignore it completely. The compromise we have adopted is to use elementary homological methods---exact sequences, diagrams, etc.---but to stop short of any results requiring a deep study of homology. In this way we hope to prepare the ground for a systematic course on homological algebra which the reader should undertake if he wishes to pursue algebraic geometry in any depth.

We have provided a substantial number of exercises at the end of each chapter. Some of them are easy and some of them are hard. Usually we have provided hints, and sometimes complete solutions, to the hard ones. We are indebted to Mr. R. Y. Sharp, who worked through them all and saved us from error more than once.

We have made no attempt to describe the contributions of the many mathematicians who have helped to develop the theory as expounded in this book. We would, however, like to put on record our indebtedness to J.-P. Serre and J. Tate from whom we learnt the subject, and whose influence was the determining factor in our choice of material and mode of presentation.

\begin{thebibliography}{9}
\bibitem{Bourbaki}
N. BOURBAKI, \emph{Alg\`ebre Commutative}, Hermann, Paris (1961--65).

\bibitem{Grothendieck}
A. GROTHENDIECK and J. DIEUDONN\'E, \emph{\'El\'ements de G\'eom\'etrie Alg\'ebrique}, Publications Math\'ematiques de l'I.H.E.S., Nos. 4, 8, 11, \ldots, Paris (1960-- ).

\bibitem{Northcott}
D. G. NORTHCOTT, \emph{Ideal Theory}, Cambridge University Press (1953).

\bibitem{ZariskiSamuel}
O. ZARISKI and P. SAMUEL, \emph{Commutative Algebra I, II}, Van Nostrand, Princeton (1958, 1960).
\end{thebibliography}


\chapter*{Notation and Terminology}
\addcontentsline{toc}{chapter}{Notation and Terminology}

Rings and modules are denoted by capital italic letters, elements of them by small italic letters. A field is often denoted by $k$. Ideals are denoted by small German characters. $\Z$, $\Q$, $\R$, $\C$ denote respectively the ring of rational integers, the field of rational numbers, the field of real numbers and the field of complex numbers.

Mappings are consistently written on the left, thus the image of an element $x$ under a mapping $f$ is written $f(x)$ and not $(x)f$. The composition of mappings $f \colon X \to Y$, $g \colon Y \to Z$ is therefore $g \circ f$, not $f \circ g$.

A mapping $f \colon X \to Y$ is \emph{injective} if $f(x_1) = f(x_2)$ implies $x_1 = x_2$; \emph{surjective} if $f(X) = Y$; \emph{bijective} if both injective and surjective.

The end of a proof (or absence of proof) is marked thus~$\blacksquare$.

Inclusion of sets is denoted by the sign $\subseteq$. We reserve the sign $\subset$ for strict inclusion. Thus $A \subset B$ means that $A$ is contained in $B$ and is not equal to $B$.
